%%=============================================================================
%% Samenvatting
%%=============================================================================

% TODO: De "abstract" of samenvatting is een kernachtige (~ 1 blz. voor een
% thesis) synthese van het document.
%
% Een goede abstract biedt een kernachtig antwoord op volgende vragen:
%
% 1. Waarover gaat de bachelorproef?
% 2. Waarom heb je er over geschreven?
% 3. Hoe heb je het onderzoek uitgevoerd?
% 4. Wat waren de resultaten? Wat blijkt uit je onderzoek?
% 5. Wat betekenen je resultaten? Wat is de relevantie voor het werkveld?
%
% Daarom bestaat een abstract uit volgende componenten:
%
% - inleiding + kaderen thema
% - probleemstelling
% - (centrale) onderzoeksvraag
% - onderzoeksdoelstelling
% - methodologie
% - resultaten (beperk tot de belangrijkste, relevant voor de onderzoeksvraag)
% - conclusies, aanbevelingen, beperkingen
%
% LET OP! Een samenvatting is GEEN voorwoord!

%%---------- Nederlandse samenvatting -----------------------------------------
%
% TODO: Als je je bachelorproef in het Engels schrijft, moet je eerst een
% Nederlandse samenvatting invoegen. Haal daarvoor onderstaande code uit
% commentaar.
% Wie zijn bachelorproef in het Nederlands schrijft, kan dit negeren, de inhoud
% wordt niet in het document ingevoegd.

\IfLanguageName{english}{%
\selectlanguage{dutch}
\chapter*{Samenvatting}
\lipsum[1-4]
\selectlanguage{english}
}{}

%%---------- Samenvatting -----------------------------------------------------
% De samenvatting in de hoofdtaal van het document

\chapter*{\IfLanguageName{dutch}{Samenvatting}{Abstract}}

IT-dienstverleners beheren voor honderden klanten gevoelige technische informatie: serveradressen, inloggegevens, netwerkconfiguraties en troubleshooting-documentatie. Bij E.S.C. BV groeide deze informatie historisch aan in Microsoft OneNote, verspreid over meerdere SharePoint-locaties per klant. Dit leidde tot structurele zoekproblemen, ontbrekende governance en aanzienlijke beveiligingsrisico's: OneNote had lange tijd geen ondersteuning voor sensitivity labels, hield niet bij wie credentials raadpleegde, en sloot daardoor niet aan bij de PAM-principes die vandaag nodig zijn. Wanneer één klant notitieboeken heeft op zowel de interne projectsite, de klantgerichte SharedWithCustomer-site als de centrale klantenhub, is tijdverlies bij het zoeken naar informatie onvermijdelijk.

E.S.C. BV stapte over naar Remote Desktop Manager (RDM) van Devolutions als centraal kennisplatform. Daarmee ontstond een concrete migratiebehoefte: hoe kan bestaande klantinformatie op een geautomatiseerde, veilige en schaalbare manier worden overgezet van OneNote naar RDM?

In dit onderzoek werden drie mogelijke extractiemethoden geëvalueerd. COM-automatisering via de OneNote Desktop-applicatie bleek onbetrouwbaar door blokkades van Windows-beveiligingsbeleid en antivirussoftware. Handmatige export is te tijdrovend bij honderden klantomgevingen. Microsoft Graph API vormt de meest haalbare aanpak, maar de vereiste \texttt{Notes.Read.All}-scope is op zichzelf breed geformuleerd. In deze implementatie werd dat correct afgebakend door uitsluitend gedelegeerde permissies te gebruiken en gebruikers enkel toegang te geven tot de relevante SharePoint-sites.

Als proof-of-concept werd een operationele batchpipeline uitgewerkt. Op orchestratieniveau gebeuren export en import via \texttt{Bulk-Export-OneNote-Queue.ps1} en \texttt{Bulk-Import-Client-ToRDM2.ps1}. In de exportketen verwerkt \texttt{Export-OneNote-Deduplicated.ps1} via Microsoft Graph API alle pagina's, secties, afbeeldingen en bijlagen van de relevante SharePoint-sites per klant, met throttling-beheer (circuit breaker, exponentiële back-off en meerfasige herstelrondes). Een deduplicatiestap groepeert pagina's op sectienaam en paginatitel en behoudt per groep de meest recente versie, zodat dubbele exports over meerdere sites niet onnodig worden geïmporteerd. In de importketen verwerkt \texttt{Import-ToRDM2.ps1} de geëxporteerde HTML-bestanden naar de juiste klantvault in RDM via de Devolutions.PowerShell-module. Daarbij wordt de tekstinhoud ook opgenomen in \emph{Description} en \emph{Documentation}, zodat inhoudszoekopdrachten bruikbaar blijven. Aanmelden gebeurt via een Entra~ID-appregistratie met eenmalige code. Voor bulktoegang wordt \texttt{Admin-Grant-SharedSites-From-CSV.ps1} gebruikt, zodat gebruikers alleen toegang krijgen tot sites die voor de export nodig zijn en de effectieve leesrechten beperkt blijven tot de toegewezen sites.

Ook op vlak van gebruiksgemak is de verbetering duidelijk. Voor de migratie moesten medewerkers eerst manueel in Teams lid worden van het juiste team (zoeken op \texttt{P-\{klantnaam\}}), daarna het SharedWithCustomer-kanaal zichtbaar maken en pas dan de juiste OneNote in de browser openen. Voor de collega's die exportbatches op de server uitvoeren, werd dit proces geautomatiseerd via de beheersscripts voor bulktoegang. Na import in RDM volstaat het om in de klantvault de map \enquote{OneNote export} te openen: de recente pagina's zijn al gededupliceerd en de informatie uit verschillende bronnotitieboeken is samengebracht in één consistente weergave.

De resultaten tonen aan dat geautomatiseerde migratie van OneNote naar RDM technisch haalbaar is. De voornaamste beperkingen liggen bij de brede formulering van \texttt{Notes.Read.All} op scope-niveau en de tijdsintensieve verwerking bij agressieve API-throttling. In de praktijk werd dit rechtenrisico beheerst via gedelegeerde permissies en gerichte sitetoekenning per gebruiker. AI-gestuurde categorisatie (zoals met Ollama) werd wel verkend als mogelijke vervolgstap, maar viel buiten de scope van deze bachelorproef.

Dit onderzoek biedt een praktische basis voor bedrijfsbrede uitrol binnen E.S.C. BV en is relevant voor IT-dienstverleners die staan voor gelijkaardige migratieuitdagingen vanuit ongestructureerde kennisplatformen naar gecentraliseerde, beveiligde informatiesystemen.
